%% kapa: exam=1 round=33 year=2022 period=1 form=A subject=economics subject_ko=경제학원론 seq=38 number=78
%% source: Q-Net 감정평가사 기출문제 (cst00309) artlSeq=5209997 file=raw/1008/5209997/2022년도 제33회 감정평가사 1차시험 1교시(민법,경제학,부동산학).pdf
%% conversion: kordoc 4.24.51 -> tools/gen1.py (auto); review=proofread-sonnet; flags=none
\begin{question}{78}
( )에 들어갈 내용으로 옳은 것은? (단, 전염병이 발생하기 전의 경제는 균형 상태이고, 총공급곡선은 우상향하고 총수요곡선은 우하향한다.)
\begin{framedtext}
폐쇄경제 AD-AS 모형에서 전염병의 발생으로 인하여 총수요와 총공급이 모두 감소할 때, 균형국민소득은 ( ㄱ )하고 균형물가수준은 ( ㄴ )하(한)다.
\end{framedtext}
\begin{choices}
  \item ㄱ: 감소, ㄴ: 감소
  \item ㄱ: 불확실, ㄴ: 불변
  \item ㄱ: 감소, ㄴ: 증가
  \item ㄱ: 불변, ㄴ: 불변
  \item ㄱ: 감소, ㄴ: 불확실
\end{choices}
\end{question}
\answer{5}
